\documentclass[11pt]{article}
\usepackage[T1]{fontenc}
\usepackage{textcomp}
\usepackage[margin=0.75in]{geometry}
\usepackage{amsmath,amssymb,amsthm}
\usepackage{array,booktabs}
\usepackage{hyperref}
\setlength{\parindent}{0pt}
\newtheorem{theorem}{Theorem}
\theoremstyle{definition}
\newtheorem{definition}{Definition}
\title{Flow Proof Artifact: parallel-interior-alternate-chain}
\author{Generated by \texttt{flow doc proof}}
\date{Flow Proof Book}
\begin{document}
\maketitle
Interior alternate angles chain to corresponding angles.\\[0.5em]
\textbf{Source.} euclid — Elements, Book I, Proposition 29\\[0.5em]
\bigskip
\noindent{\large \textbf{Derived fact 1.} \textit{interior alternate alpha equals corresponding beta} \hfill the Euclidean plane \cdot parallel lines \cdot interior alternate chains to corresponding}\par
\smallskip\noindent\textit{Coordinate: the Euclidean plane \cdot parallel lines \cdot interior alternate chains to corresponding}\par
\smallskip\noindent\textit{Source: euclid}\par
\smallskip\noindent\textit{Built on: corresponding and f angles chain, for parallel lines in on the Euclidean plane, alternate and z angles chain, for parallel lines in on the Euclidean plane}\par
\medskip
\noindent\textbf{Goal.} interior alternate alpha equals corresponding beta\par
\noindent$\displaystyle \angle \text{interior alternate alpha} = \angle \text{corresponding beta}$\par
\medskip
\noindent
\begin{tabular}{@{} >{\raggedright\arraybackslash}p{0.06\textwidth} >{\raggedright\arraybackslash}p{0.47\textwidth} >{\raggedleft\arraybackslash}p{0.41\textwidth} @{}}
\textbf{\#} & \textbf{Proof} & \textbf{Mathematics} \\
\midrule
\textbf{1.} & We prove that interior alternate chains to corresponding for parallel lines in the Euclidean plane. & \\[0.45em]
\textbf{2.} & We invoke the derived fact governing parallel lines in the Euclidean plane: corresponding and f angles chain, for parallel lines in on the Euclidean plane. & \\[0.45em]
\textbf{3.} & We invoke the derived fact governing parallel lines in the Euclidean plane: alternate and z angles chain, for parallel lines in on the Euclidean plane. & \\[0.45em]
\textbf{4.} & From step 2 and step 3, this implies angle interior alternate alpha equals angle corresponding beta. Hence proven. & $\displaystyle \angle \text{interior alternate alpha} = \angle \text{corresponding beta}$ \\[0.45em]
\bottomrule
\end{tabular}
\label{thm:Geometry-parallel lines-interior_alternate_chains_to_corresponding}
\par\medskip\hrule
\end{document}
